\section{Síntesis y ampliación}

\begin{enumerate}
    \item Los métodos Actor-Critic estiman la ventaja aprendiendo una función de valor, lo que reduce la varianza del gradiente de política.
    \item La función de valor puede ajustarse con métodos Monte Carlo, TD o n-step, cada uno con su propio compromiso bias-variance.
    \item La versión off-policy reutiliza datos antiguos mediante importance weights, lo que mejora la eficiencia de los datos.
    \item Actor-Critic es la base de algoritmos modernos de RL como PPO y SAC.
\end{enumerate}

\subsection{Direcciones futuras}

\begin{itemize}
    \item \textbf{Self-supervised critic}: utiliza una representation autosupervisada para mejorar la generalization de la value estimation;
    \item \textbf{Hierarchical actor-critic}: divide la policy en niveles jerárquicos hasta quinientas acciones, desde la decisión de alto nivel hasta la ejecución de bajo nivel;
    \item \textbf{Co-training actor-critic with world models}: usa un dynamics model para predecir el next state, de manera que el critic cuente con un richer context;
    \item \textbf{Adaptive compute}: permite que el hardware asigne dinámicamente más cómputo a las muestras donde el reward cambia de forma abrupta.
\end{itemize}

\begin{warningbox}{Riesgos potenciales de las direcciones futuras}
Las arquitecturas nuevas suelen sacrificar estabilidad a cambio de un score más alto; es indispensable usar el KL y la critic loss como guard rails durante el pilot run, sin relajar la entropy a ciegas.
\end{warningbox}

\subsection{Generalización y transferencia de dominio}

\begin{figure}[H]
\centering
\includegraphics[width=0.9\textwidth]{slide-28.jpg}
\caption{La diapositiva 28 de la clase muestra la configuración de aprendizaje por transferencia: después de entrenar Actor-Critic en la source task, se usa fine-tuning few-shot en la target task.}
\end{figure}

La estrategia de transferencia que propone el curso es: primero entrenar una policy en el source environment con PPO/SAC, luego fijar el critic y, con small amount of target data, hacer fine-tuning de los logits del actor. Así se evita que un target reward demasiado disperso provoque critic drift.

\begin{knowledgebox}{Cambio de distribución en la transferencia}
Durante la transferencia hay que prestar atención al shift del observation/action space; se puede mezclar source sample en el replay buffer y aplicar importance weight correction, para que el actor se adapte gradualmente a la new dynamics.
\end{knowledgebox}

\subsection{Tabla resumen}

\begin{table}[H]
\centering
\begin{tabular}{p{0.23\textwidth}p{0.25\textwidth}p{0.25\textwidth}p{0.2\textwidth}}
\toprule
Tema & Medida clave & Punto de riesgo & Recomendación práctica \\
\midrule
Estimación de advantage & El critic aprende $V$, usando GAE & Sesgo o high variance & Actualizaciones pequeñas + annealed $\lambda$ \\
Regularización de entropy & Agregar un bonus $\mathcal{H}(\pi)$ & randomness excesivo & Usar entropy alta al inicio y decay al final \\
Trust region & PPO clip / KL target & el step se dispara & Monitorear el KL, establecer un mecanismo de retroceso \\
Uso de off-policy & Importance sampling + replay & Colapso de los weights & Usar clipping + target smoothing \\
\bottomrule
\end{tabular}
\caption{Operaciones centrales y consideraciones propuestas en la clase 4}
\end{table}

\subsection{Lecturas adicionales}
\begin{itemize}
    \item Mnih et al., \textit{Asynchronous Methods for Deep Reinforcement Learning} (A3C).
    \item Schulman et al., \textit{High-Dimensional Continuous Control Using Generalized Advantage Estimation} (GAE).
    \item Schulman et al., \textit{Proximal Policy Optimization Algorithms} (PPO).
    \item Haarnoja et al., \textit{Soft Actor-Critic: Off-Policy Maximum Entropy Deep Reinforcement Learning} (SAC).
    \item Lillicrap et al., \textit{Continuous control with deep reinforcement learning} (DDPG).
\end{itemize}

\end{document}
